\documentclass[11pt]{article}
\usepackage[margin=1in]{geometry}
\usepackage[T1]{fontenc}
\usepackage[utf8]{inputenc}
\usepackage{amsmath,amssymb,mathtools}
\usepackage{booktabs}
\usepackage{array}
\usepackage{graphicx}
\usepackage{hyperref}
\usepackage{xcolor}

\title{Simultaneous Multitone SQR in Dispersive cQED:\\
When One Common Waveform Can and Cannot Drive Multiple Branches}
\author{}
\date{\today}

\begin{document}
\maketitle

\section{Motivation}

The earlier SQR waveform study showed that strict logical performance can fail
even when relaxed branchwise metrics look strong. This follow-up asks a harder
question: can one common multitone waveform implement \emph{multiple} selective
qubit rotations at the same time? If the answer is no for the simplest family,
we also want to know whether the failure is primarily:
\begin{enumerate}
    \item a removable cavity block-phase profile,
    \item a simple pulse-calibration problem,
    \item or a deeper limitation of the common Gaussian multitone ansatz.
\end{enumerate}

The key scientific caution in this study is that fidelity alone can be
misleading for small target angles. If a pulse acts almost like identity, then
it can still score well against a small-angle target simply because the target
itself is close to identity. We therefore benchmark both fidelity and the
actually executed target-branch rotation angle.

\section{System \& Methods}

We model a dispersive transmon-cavity system with
\[
\omega_q = 2\pi \times 6.150~\mathrm{GHz}, \qquad
\omega_c = 2\pi \times 5.241~\mathrm{GHz}, \qquad
\alpha = 2\pi \times (-255~\mathrm{MHz}), \qquad
\chi = 2\pi \times (-2.84~\mathrm{MHz}).
\]
The logical cavity window is \(n = 0,1,2,3\). The main targeted-subspace study
uses \(n_{\mathrm{tr}}=2\) and \(n_{\mathrm{cav}}=6\), because the current
`cqed\_sim` targeted-subspace multitone workflow evaluates conditioned qubit
metrics through a two-level reduced qubit state. Representative waveforms are
then replayed on \(n_{\mathrm{tr}}=3\) to estimate \(|f\rangle\) leakage.

We study four simultaneous target sets:
\[
\{0,1\}, \qquad \{0,2\}, \qquad \{0,1,2\}, \qquad \{0,1,2,3\},
\]
with target rotation angles
\[
\theta/\pi \in \{1/8,\, 1/4,\, 1/2,\, 1\},
\]
always around the \(X\) axis. The duration scan uses
\[
\chi T / 2\pi \in \{1,2,3,5\}.
\]

For each point we compute:
\begin{itemize}
    \item reduced conditioned multitone fidelity,
    \item strict logical targeted-subspace fidelity,
    \item best-fit cavity block-phase compiled fidelity,
    \item state-transfer fidelity,
    \item same-block population preservation,
    \item best-fit cavity block phases,
    \item and the mean simulated target-branch Bloch angle \(\bar{\theta}_{\rm sim}\).
\end{itemize}

We additionally test three correction ideas on a representative hard case
(pair target \(\{0,1\}\), \(\theta=\pi\), \(\chi T/2\pi = 3\)):
\begin{enumerate}
    \item common target-tone amplitude correction \(\Delta \lambda\),
    \item multistart amplitude/phase/detuning optimization,
    \item naive repeated-step compilation using eight compiled \(\pi/8\) steps.
\end{enumerate}

Validation includes sanity checks, convergence checks in \(dt\) and
\(n_{\mathrm{cav}}\), and qutrit replay spot-checks.

\section{Results}

\subsection{Baseline Fidelity Maps Are Misleading Without an Angle-Response Diagnostic}

Figure~\ref{fig:conditioned} shows the reduced conditioned fidelity over the full
scan. Taken at face value, the plots could suggest that simultaneous multitone
SQR is good for small angles and only breaks down for large angles.

\begin{figure}[ht]
    \centering
    \includegraphics[width=\textwidth]{../figures/simultaneous_multitone_conditioned_heatmaps.png}
    \caption{Reduced conditioned fidelity for simultaneous multitone targets.
    Fidelity decreases with requested angle, but this plot alone does not tell us
    whether the pulse is really implementing the requested rotation or merely
    staying close to identity.}
    \label{fig:conditioned}
\end{figure}

The crucial correction is Figure~\ref{fig:angle-response}. At the representative
\(\chi T/2\pi = 3\) slice, the mean simulated target angle is tiny:
\[
\bar{\theta}_{\rm sim}/\pi \sim 10^{-5}\text{--}10^{-4},
\]
even when the requested target is \(\theta=\pi\). For the pair-adjacent case
\(\{0,1\}\), the target-angle response ratios are
\[
\bar{\theta}_{\rm sim}/\theta \approx
2.23\times 10^{-4},\;
2.23\times 10^{-4},\;
2.22\times 10^{-4},\;
2.20\times 10^{-4}
\]
for \(\theta/\pi = 1/8, 1/4, 1/2, 1\), respectively. The other target sets show
the same behavior at the \(10^{-4}\) level.

\begin{figure}[ht]
    \centering
    \includegraphics[width=0.74\textwidth]{../figures/simultaneous_multitone_target_angle_response.png}
    \caption{Actual target-angle response at \(\chi T/2\pi = 3\). The ideal
    diagonal line would mean the waveform implements the requested rotation. The
    simulated response instead stays near zero for every simultaneous target set,
    so the nominal ``good'' small-angle fidelity is mostly trivial identity
    overlap.}
    \label{fig:angle-response}
\end{figure}

This changes the interpretation completely: the common Gaussian simultaneous
multitone family is not a successful multi-SQR pulse that merely struggles at
large angle. It is, to leading order, an almost-identity pulse on every block.
The apparently strong small-angle performance arises because \(R_X(\theta)\) is
itself close to identity for small \(\theta\).

\subsection{Strict Logical Validation Shows Block Preservation Without Useful Rotation}

Figure~\ref{fig:strict} gives the strict targeted-subspace logical fidelity.
The scan is nearly independent of \(\chi T/2\pi\) over the explored range:
changing duration from \(1\) to \(5\) improves the representative pair-target
\(\theta=\pi\) strict fidelity only from \(0.2476\) to \(0.2483\).

\begin{figure}[ht]
    \centering
    \includegraphics[width=\textwidth]{../figures/simultaneous_multitone_strict_heatmaps.png}
    \caption{Strict logical targeted-subspace fidelity. These values are lower
    than the reduced conditioned metrics, but the same qualitative warning
    remains: small-angle numbers are dominated by identity overlap, while
    large-angle simultaneous targets fail strongly.}
    \label{fig:strict}
\end{figure}

Representative \(\chi T/2\pi=3\) strict fidelities are:
\begin{center}
\begin{tabular}{>{\raggedright\arraybackslash}p{3.0cm}cccc}
\toprule
Target set & \(\theta/\pi=1/8\) & \(1/4\) & \(1/2\) & \(1\) \\
\midrule
\(\{0,1\}\) & 0.9759 & 0.9205 & 0.7242 & 0.2477 \\
\(\{0,2\}\) & 0.9759 & 0.9205 & 0.7244 & 0.2479 \\
\(\{0,1,2\}\) & 0.9664 & 0.8843 & 0.6051 & 0.0616 \\
\(\{0,1,2,3\}\) & 0.9570 & 0.8490 & 0.4970 & \(3.8\times 10^{-7}\) \\
\bottomrule
\end{tabular}
\end{center}

These values look respectable for small \(\theta\), but Figure~\ref{fig:angle-response}
shows that they do \emph{not} correspond to correct simultaneous rotations.

At the same time, the pulse preserves cavity blocks extremely well: the minimum
same-block population mean over the whole scan is
\[
0.999998207.
\]
Representative qutrit replay spot-checks also show tiny transmon leakage, with a
maximum observed \(|f\rangle\) population of only
\[
7.08\times 10^{-9}.
\]
So the dominant failure is neither block mixing nor obvious qutrit leakage. The
waveform simply fails to generate the requested qubit rotations.

\subsection{Cavity Block-Phase Compilation Is Real but Secondary}

The best-fit cavity block-phase profile is extremely regular. At
\(\chi T/2\pi = 3\), a representative best-fit phase vector is approximately
\[
\phi_n \approx (0,\,-0.03267,\,-0.06535,\,-0.09802)\ {\rm rad},
\]
which is nearly linear in \(n\). However, the compiled gain remains tiny:
the largest best-fit cavity-compilation gain anywhere in the scan is only
\[
\Delta F_{\rm comp}^{\max} = 1.47\times 10^{-3}.
\]

\begin{figure}[ht]
    \centering
    \includegraphics[width=\textwidth]{../figures/simultaneous_multitone_compiled_gain_heatmaps.png}
    \caption{Best-fit cavity-only compilation gain. The induced block phase is
    smooth and real, but correcting it changes the conclusion only at the
    \(10^{-3}\) level. This is not primarily a cavity-phase-compilation problem.}
    \label{fig:gain}
\end{figure}

So a cavity phase exists, but it is not the main obstacle. The main obstacle is
that the common waveform does not perform the intended block-local rotation in
the first place.

\subsection{Simple Pulse Corrections Do Not Rescue the Hard Regime}

Figure~\ref{fig:corrections} summarizes the correction studies for the hard
pair-target \(\pi\)-gate at \(\chi T/2\pi=3\).

\begin{figure}[ht]
    \centering
    \includegraphics[width=\textwidth]{../figures/simultaneous_multitone_correction_checks.png}
    \caption{Correction checks for the representative hard case. Left:
    sweeping a common target-tone amplitude correction is essentially flat.
    Middle: broad multistart amplitude/phase/detuning correction remains stuck
    near the same reduced fidelity. Right: qutrit replay spot-checks confirm the
    failure is not driven by large \(|f\rangle\) leakage.}
    \label{fig:corrections}
\end{figure}

The explicit results are:
\begin{itemize}
    \item Sweeping a common target-tone amplitude correction
    \(\Delta \lambda \in [-1.5, 3.0]\) changes the reduced mean fidelity only
    from \(0.50000000\) to \(0.50000007\).
    \item Three multistart waveform-side optimizations over amplitude, phase,
    and detuning reach a best reduced fidelity of only \(0.50000040\). The
    best corresponding strict logical fidelity is \(0.231997\), actually below
    the raw baseline strict fidelity \(0.247713\).
    \item Thus, simple pulse correction within this ansatz does not rescue the
    simultaneous \(\pi\)-gate regime.
\end{itemize}

\subsection{Naive Repeated Small-Angle Compilation Also Fails}

One possible fallback idea is to use many high-fidelity small-angle primitives
instead of one direct \(\pi\) pulse. We tested this by composing the best-fit
cavity-compiled restricted operator for the \(\pi/8\) step eight times and
comparing the resulting logical operator with the desired direct \(\pi\) target.

\begin{figure}[ht]
    \centering
    \includegraphics[width=0.72\textwidth]{../figures/simultaneous_multitone_segmented_vs_direct.png}
    \caption{Naive repeated-step compilation is not enough. Eight compiled
    \(\pi/8\) steps do not outperform the already-bad direct \(\pi\) gate for
    either the pair or triple target sets.}
    \label{fig:segmented}
\end{figure}

The resulting operator-level fidelities are:
\[
F_{\rm seg}^{\{0,1\}} = 0.1546, \qquad
F_{\rm seg}^{\{0,1,2\}} = 0.0311,
\]
compared with direct-\(\pi\) strict fidelities \(0.2477\) and \(0.0616\),
respectively. Naive repeated small steps therefore do not solve the problem.

\section{Discussion}

The central conclusion is that the common Gaussian simultaneous multitone family
studied here is \emph{not} a practical simultaneous multi-SQR primitive. The
strongest evidence is not the fidelity maps alone but the explicit target-angle
response: the waveform leaves every target branch almost unrotated, with only a
tiny \(O(10^{-4})\) fractional response to the requested angle.

This also explains several otherwise puzzling features:
\begin{itemize}
    \item Pair-adjacent and pair-separated targets perform almost identically,
    so nearest-neighbor branch crowding is not the main issue.
    \item Increasing \(\chi T/2\pi\) over the scanned range barely changes the
    outcome, so duration alone does not fix the ansatz.
    \item The cavity block phase is smooth and nearly linear, but correcting it
    helps only at the \(10^{-3}\) level because the principal defect is missing
    rotation, not missing phase compilation.
    \item Tiny qutrit leakage and excellent same-block preservation show that
    the pulse is coherent and block-preserving; it is just the wrong coherent
    action.
\end{itemize}

From an experiment-facing control perspective, the lesson is very specific:
before trusting any simultaneous-target fidelity metric, one must calibrate the
actual target-branch angle response. In this study, the nominal parameter
\(\theta\) in the simultaneous common-waveform builder is not a reliable proxy
for the executed branch rotation.

The likely next steps are therefore not more minor corrections to this exact
family, but richer waveform classes:
\begin{enumerate}
    \item segmented simultaneous pulses with independently shaped subwindows,
    \item explicitly time-dependent tone phases or chirps,
    \item or direct optimal control constrained by hardware bandwidth.
\end{enumerate}
Only after such a waveform already produces the correct block-local rotations
does cavity phase compilation become a meaningful secondary cleanup layer.

\section{Conclusion}

The study answers the original question in a fairly decisive way.

\begin{center}
\begin{tabular}{>{\raggedright\arraybackslash}p{4.0cm}p{2.2cm}p{2.4cm}p{5.0cm}}
\toprule
Primitive idea & Verdict & Main issue & Recommendation \\
\midrule
Direct common Gaussian simultaneous multitone SQR & Not viable & Target branches barely rotate; fidelity can look good only because small-angle targets are near identity & Do not deploy as a simultaneous multi-SQR gate \\
Direct pulse + cavity block-phase compilation & Not enough & Compilation gain is only \(O(10^{-3})\) & Use only after a better waveform family is found \\
Simple amplitude correction & Not enough & Response is essentially flat in the hard regime & Not worth pursuing alone \\
Broad multistart amplitude/phase/detuning correction & Not enough & No meaningful gain over baseline & Treat as evidence of ansatz limitation, not just poor initialization \\
Naive repeated small-angle compilation & Not enough & Composition of weak near-identity steps does not build the desired \(\pi\) gate & Do not use without a richer inter-step compilation strategy \\
\bottomrule
\end{tabular}
\end{center}

The best immediate recommendation is:
\begin{quote}
Do \emph{not} use the common Gaussian simultaneous multitone family as an
experiment-facing simultaneous multi-SQR primitive. Move instead to richer
segmented or optimal-control waveforms, and evaluate them with explicit
target-angle response diagnostics in addition to fidelity.
\end{quote}

The most important unresolved bottleneck is not cavity phase. It is the absence
of a waveform family that can produce the desired simultaneous block-local
rotations at all.

\end{document}
